\scanpage{342}
% @ top centre
\mspage{66}

% @ top, displayed question
\[
 F\overset{?}{\Longrightarrow} A,
\]
% @ upper left beside question
% slanted
\margin{C'est faux, cf plus bas l'implication $\unsure{\mathrm{Doc}|\mathrm{Loc}}(7D_{W_\infty})\Longrightarrow \mathrm{Loc}(7\mathrm F)$ est \uline{stricte}.}
% @ top paragraph below question
du moins \unsure{ss|sous} des conditions convenables assez anodines sur $W$, genre $\mathrm{Loc}(1\ \mathrm{ter},3\ \mathrm{bis},8\ \mathrm{bis})$.

% @ upper to middle paragraph
C'est sans doute là la question la plus substantielle \inserted{et la plus intéressante} qui est soulevée par le diagramme de la page 63, et qui ne fait plus appel à l'introduction d'un localiseur fondamental $W_0$ auxiliaire. \deleted{Après}\inserted{Vu} le lemme de Brown, il ne me semble pas absolument exclu qu'un tel miracle ait bel et bien lieu !

% @ centre, interlined heading and section opening
4.4.15. \inserted{(Digression sur l'axiome $\mathrm{Loc}(6)$ des cofibrations)} Cela attire mon attention sur la question en quelque sorte duale : Si $\mathrm{Cat}$ a assez de $W$-\uline{co}fibrations, i.e si \unsure{muni|munie} de $(W,\mathrm{Cof}_W)$ c'est une catégorie à cofibrations, cela implique-t-il déjà l'axiome $\mathrm{Loc}(6\mathrm a)$\unsure{$*$|.|$?$} (qui avait été dégagé, essentiellement, pour assurer qu'il en soit bien ainsi). La proposition 3.6. \inserted{(p.25)} nous dit qu'il revient au \unsure{$=$|même} de prouver que, si $W$ satisfait $\mathrm{Cof}$, alors les carrés de M.V. ouverts sont
% archive sheet number: 341
